\documentclass[10pt,a4paper]{amsart}
\usepackage[margin=19mm]{geometry}
\usepackage{amsmath,amssymb,mathtools}
\usepackage[scr=rsfs,cal=euler]{mathalfa}
\usepackage{xfrac}
\usepackage[unicode,hidelinks,bookmarks=false]{hyperref}
\newcommand{\ie}{i.e.\ }
\newcommand{\cf}{cf.\ }
\newcommand{\resp}{respectively\ }
\newcommand{\Subsection}{\subsection}
\providecommand{\nobreakdash}{\nobreak\hbox{-}}
\setlength{\emergencystretch}{2em}
\multlinegap0pt
\usepackage{bm}
\usepackage{bbm}
\usepackage{dsfont}
\usepackage{xspace}
\usepackage{cancel}
\usepackage{mathtools}
\usepackage{amsthm}
\makeatletter
\@ifundefined{theorem}{\newtheorem{theorem}{Theorem}}{}
\@ifundefined{lemma}{\newtheorem{lemma}[theorem]{Lemma}}{}
\makeatother
\makeatletter
\newcommand{\C}{\mathbb C}
\newcommand{\cB}{\mathcal B}
\DeclareMathOperator{\ord}{ord}
\expandafter\ifx\csname poc\endcsname\relax
\newenvironment{poc}
{\begin{proof}[Proof of the claim]\renewcommand*{\qedsymbol}{$\blacksquare$}}
\fi
\makeatother
\title[When Do Subset Sums in Finite Abelian Groups Support $2$-Des]{When Do Subset Sums in Finite Abelian Groups Support $2$-Designs?}
\author{Hengfeng Liu, Chunming Tang, Cuiling Fan, Zhengchun Zhou}
\date{Source: arXiv:2607.24426v1 (CC BY 4.0)}
\begin{document}
\maketitle

\noindent\emph{Source and licence.} When Do Subset Sums in Finite Abelian Groups Support $2$-Designs?. Version arXiv:2607.24426v1 (CC BY 4.0), distributed under the Creative Commons Attribution 4.0 International licence, as recorded in the arXiv licence metadata for this exact version and shown on the abstract page https://arxiv.org/abs/2607.24426v1. Full rights notice: licenses/arxiv-2607.24426-CC-BY-4.0.txt.\par\smallskip

\noindent\textit{Editorial scope.} Counted unit: the Theorem whose source label is \texttt{thm:energy-expansion} in the article (source0.tex lines 717--732, source label \texttt{thm:energy-expansion}), delivered together with the article's own complete original proof of it (source0.tex lines 734--919, one \texttt{proof} environment of 186 lines). No mathematics is written, repaired or re-derived: statement and proof are the article's own text line for line. Required context retained uncounted: lem:product-identity (lines 224--232). Every definition, display and label of those lines is retained unchanged. Omitted from this package and not redistributed: 1--223, 233--716, 733--733, 920--1536. Nothing inside a retained excerpt is omitted or altered: every retained body is delivered exactly as the article's lines are written, so no omission receipt of any kind accompanies this package. Citations: the retained excerpts carry no citation command at all, so no bibliography entry of the article is transcribed and no reference is redistributed. Reference adaptation: none was needed and none was made; every reference inside the delivered unit resolves inside it, so no unresolved reference or citation is produced. Printed slips kept literal: no printed word, symbol, hypothesis or number of the counted statement or of its proof was altered. Layout and class: the article's own class is \texttt{article} with packages including amsthm, amsmath, amssymb, amsfonts, url, booktabs, tikz, setspace, fancyhdr, bm, geometry, hyperref, enumerate, enumitem; the wrapper is the lane's pinned \texttt{amsart} editorial wrapper at 10pt on A4, which restates the theorem-like environments the retained text needs and adds the article's own macro definitions that the retained text needs (\texttt{\textbackslash C}, \texttt{\textbackslash cB}, \texttt{\textbackslash ord}), with extra wrapper packages (bm, bbm, dsfont, xspace, cancel, mathtools). The delivered numbering is the unit's own. Assets: no retained span contains a figure environment, a graphics-inclusion command, a PDF-page inclusion, a table or a TikZ picture; no image file is copied into this package, the package declares no asset dependency and no image file is redistributed with it. Licence: version arXiv:2607.24426v1 of this article is distributed under the Creative Commons Attribution 4.0 International licence, as recorded in the arXiv licence metadata for this exact version and shown on the abstract page https://arxiv.org/abs/2607.24426v1. A full rights notice is delivered at licenses/arxiv-2607.24426-CC-BY-4.0.txt. Characters: every retained excerpt is pure ASCII, so the delivered engine, XeLaTeX with its default Latin Modern text font, reports no missing character.
\begin{lemma}
	\label{lem:product-identity}
	Let \(\chi\in\widehat G\) have order \(d\).  Then
	\[
	\prod_{g\in G}\bigl(1-\chi(g)Z\bigr)
	=
	\bigl(1-Z^d\bigr)^{n/d}.
	\]
\end{lemma}

\begin{theorem}\label{thm:energy-expansion}
	There exists a constant \(\beta\), independent of \(\chi\), such that for every nontrivial \(\chi\in\widehat G\),
	\begin{equation}
	E_\chi
	=
	\beta+
	\sum_{\substack{C\le\widehat G\ \mathrm{cyclic}\\
		\langle\chi\rangle\le C}}
	A_{|C|}R_C(x),
	\end{equation}
	where $R_C(x)$ is the Ramanujan sum defined above. Consequently, if $(G,\cB_k^x)$ is a $2$-design, then
	\begin{equation}\label{eq:energy-system}
	\sum_{\substack{C\le\widehat G\ \mathrm{cyclic}\\C_0\le C}}A_{|C|}R_C(x)
	\end{equation}
	is independent of the choice of nontrivial cyclic subgroup $C_0\le\widehat G$.
\end{theorem}

\begin{proof}
	By the definition of \(E_\chi\) and character orthogonality,
	\[
	\begin{aligned}
	E_\chi
	&=
	\frac1n
	\sum_{\eta\in\widehat G}
	\eta(-x)
	\operatorname{Coef}_{k}\left\{
	\sum_{B\subseteq G}
	\left|\sum_{b\in B}\chi(b)\right|^2
	\eta\left(\sum_{b\in B}b\right)
	Z^{|B|}\right\}.
	\end{aligned}
	\]
Fix \(\eta\in\widehat G\). We rewrite the inner subset sum using the ordered pair \((u,v)\) that appears after expanding the square. Namely,
$\left|\sum_{b\in B}\chi(b)\right|^2=\sum_{u,v\in B}\chi(u)\chi(-v)$.

Hence
\begin{equation}\label{eq:pair-counting-view}
	\begin{aligned}
		&\sum_{B\subseteq G}
		\left|\sum_{b\in B}\chi(b)\right|^2
		\eta\left(\sum_{b\in B}b\right)Z^{|B|} \\
		&\qquad =
		\sum_{u,v\in G}\chi(u)\chi(-v)
		\sum_{\substack{B\subseteq G\\ u,v\in B}}
		\eta\left(\sum_{b\in B}b\right)Z^{|B|}.
	\end{aligned}
\end{equation}


Put $S_\eta(Z)=\prod_{g\in G}(1+\eta(g)Z)$ and $r_g=\frac{\eta(g)Z}{1+\eta(g)Z}$. If \(u\ne v\), then forcing \(u\) and \(v\) to lie in \(B\) gives
\begin{equation}\label{eq:off-diagonal-forced-choice}
	\begin{aligned}
		\sum_{\substack{B\subseteq G\\ u,v\in B}}
		\eta\left(\sum_{b\in B}b\right)Z^{|B|}
		&=
		\eta(u)\eta(v)Z^2
		\prod_{g\in G\setminus\{u,v\}}(1+\eta(g)Z)  \\
		&=
		S_\eta(Z)r_ur_v.
	\end{aligned}
\end{equation}
If \(u=v\), then 
\begin{equation}\label{eq:diagonal-forced-choice}
	\begin{aligned}
		\sum_{\substack{B\subseteq G\\ u\in B}}
		\eta\left(\sum_{b\in B}b\right)Z^{|B|}
		&=
		\eta(u)Z
		\prod_{g\in G\setminus\{u\}}(1+\eta(g)Z)  \\
		&=
		S_\eta(Z)r_u.
	\end{aligned}
\end{equation}
Substituting Eqs.~\eqref{eq:off-diagonal-forced-choice} and~\eqref{eq:diagonal-forced-choice} into Eq.~\eqref{eq:pair-counting-view}, we obtain
\begin{equation}\label{eq:inner-sum-offdiag-diag}
	\begin{aligned}
		&\sum_{B\subseteq G}
		\left|\sum_{b\in B}\chi(b)\right|^2
		\eta\left(\sum_{b\in B}b\right)Z^{|B|} \\
		&\qquad =
		S_\eta(Z)
		\left(\sum_{\substack{u,v\in G\\ u\ne v}}
		\chi(u)\chi(-v)r_ur_v
		+
		\sum_{u\in G}r_u
		\right).
	\end{aligned}
\end{equation}
For the first summand in Eq.~\eqref{eq:inner-sum-offdiag-diag}, we have
\begin{equation}\label{eq:offdiag-completed}
	\begin{aligned}
		\sum_{\substack{u,v\in G\\ u\ne v}}
		\chi(u)\chi(-v)r_ur_v
		&=
		\sum_{u,v\in G}\chi(u)\chi(-v)r_ur_v
		-
		\sum_{g\in G}\chi(g)\chi(-g)r_g^2  \\
		&=
	\sum_{g\in G}\chi(g)r_g
	\sum_{g\in G}\chi(-g)r_g
		-
		\sum_{g\in G}r_g^2.
	\end{aligned}
\end{equation}
Combining Eqs.~\eqref{eq:inner-sum-offdiag-diag} and~\eqref{eq:offdiag-completed}, we get
\begin{equation}\label{eq:inner-sum-final}
	\begin{aligned}
		&\sum_{B\subseteq G}
		\left|\sum_{b\in B}\chi(b)\right|^2
		\eta\left(\sum_{b\in B}b\right)Z^{|B|} \\
		&\qquad =
		S_\eta(Z)
		\left[
		\left(\sum_{g\in G}\chi(g)r_g\right)
		\left(\sum_{g\in G}\chi(-g)r_g\right)
		+
		\sum_{g\in G}(r_g-r_g^2)
		\right].
	\end{aligned}
	\end{equation}

Observe that the second term inside the brackets is independent of \(\chi\). After summing over \(\eta\) and extracting the coefficient of \(Z^k\), all such contributions are absorbed into the constant \(\beta\).

It remains to compute
\[
S_\eta(Z)
\left(\sum_{g\in G}\chi(g)r_g\right)
\left(\sum_{g\in G}\chi(-g)r_g\right).
\]
Expanding \(r_g\) as a power series, we obtain $r_g=\sum_{j\ge1}(-1)^{j-1}\eta(g)^j Z^j$.
Then we have
\[
\sum_{g\in G}\chi(g)r_g
=
\sum_{j\ge1}(-1)^{j-1}Z^j
\sum_{g\in G}(\chi\eta^j)(g).
\]
Thus the inner sum is nonzero precisely when \(\eta^j=\chi^{-1}\).
Similarly,
\[
\sum_{g\in G}\chi(-g)r_g
=
\sum_{j\ge1}(-1)^{j-1}Z^j
\sum_{g\in G}(\chi^{-1}\eta^j)(g),
\]
and here the inner sum is nonzero precisely when \(\eta^j=\chi\).
Hence the product is zero unless \(\chi\in\langle\eta\rangle\).

Assume now that \(\chi\in\langle\eta\rangle\), and let
\(d=\ord(\eta)\). Choose \(1\le a\le d-1\) such that
\(\eta^a=\chi^{-1}\). Then
\[
\sum_{g\in G}\chi(g)r_g
=
n\sum_{t\ge0}(-1)^{a+td-1}Z^{a+td}
=
n\frac{(-1)^{a-1}Z^a}{1-(-1)^d Z^d},
\]
and
\[
\sum_{g\in G}\chi(-g)r_g
=
n\sum_{t\ge0}(-1)^{d-a+td-1}Z^{d-a+td}
=
n\frac{(-1)^{d-a-1}Z^{d-a}}{1-(-1)^d Z^d}.
\]
Their product is therefore
\begin{equation}\label{eq:dependent-energy-term}
\sum_{g\in G}\chi(g)r_g
\sum_{g\in G}\chi(-g)r_g=n^2(-1)^{d-2}\frac{Z^d}{(1-(-1)^d Z^d)^2}.
\end{equation}

By Lemma~\ref{lem:product-identity} and Eq.~\eqref{eq:dependent-energy-term}, we have
\[
S_\eta(Z)
\left(\sum_{g\in G}\chi(g)r_g\right)
\left(\sum_{g\in G}\chi(-g)r_g\right)=n^2(-1)^{d-2}Z^d
\bigl(1-(-1)^d Z^d\bigr)^{n/d-2}.
\]
Therefore the contribution of this \(\eta\) to the \(\chi\)-dependent
part, including the outer factor \(n^{-1}\eta(-x)\), is
\[
\begin{aligned}
	&\frac{1}{n}\eta(-x)\operatorname{Coef}_{k}\left\{S_\eta(Z)
	\left(\sum_{g\in G}\chi(g)r_g\right)
	\left(\sum_{g\in G}\chi(-g)r_g\right)\right\} \\
	&\qquad
	=\eta(-x)\,n(-1)^{d-2}
	\operatorname{Coef}_{k-d}\left\{
	\bigl(1-(-1)^dZ^d\bigr)^{n/d-2}\right\} \\
	&\qquad
	=\eta(-x)A_d.
\end{aligned}
\]
Consequently, the $\chi$-dependent part is obtained by summing $\eta(-x)A_{\ord(\eta)}$ over all $\eta$ with $\chi\in\langle\eta\rangle$. Grouping these characters by the cyclic subgroup $C=\langle\eta\rangle$ gives
\[
\sum_{\substack{C\le\widehat G\ \mathrm{cyclic}\\
		\langle\chi\rangle\le C}}
A_{|C|}R_C(x).
\]
This completes the proof.
\end{proof}

\end{document}
